\chapter{การพัฒนาระบบฐานข้อมูลและแดชบอร์ดด้วยภาษาไพทอน}
\label{chap:dashboard}

แม้ว่าบริการคลาวด์สาธารณะอย่าง Google Firebase\index{Firebase} จะมีความสะดวกในการเริ่มต้น แต่ในบริบทของงานวิจัยวิชาการและระบบเกษตรกรรมระดับพาณิชย์ การสร้างฐานข้อมูลและระบบแดชบอร์ดขึ้นมาเอง (Self-Hosted Architecture) มอบข้อได้เปรียบเชิงยุทธศาสตร์ที่เหนือกว่า ทั้งในมิติด้านความเป็นเจ้าของข้อมูลอย่างสมบูรณ์ (Data Sovereignty) ความเร็วในการเข้าถึงข้อมูลโดยไม่มีค่าบริการแอบแฝง และความยืดหยุ่นในการบูรณาการโมเดลปัญญาประดิษฐ์ระดับลึก บทนี้นำเสนอการพัฒนาแพลตฟอร์ม Python Full-Stack ซึ่งประกอบด้วย \textbf{FastAPI\index{FastAPI}, SQLite\index{SQLite}/PostgreSQL และ Streamlit\index{Streamlit}}

\section{สถาปัตยกรรมระบบเซิร์ฟเวอร์ท้องถิ่นและความเป็นอิสระของข้อมูล}

ระบบ Self-Hosted ถูกออกแบบให้สามารถรันบนเครื่องคอมพิวเตอร์ส่วนบุคคล (Local Workstation) มินิพีซี (Mini PC) หรือเซิร์ฟเวอร์เสมือนส่วนตัว (Cloud VPS) โดยมีโครงสร้างการทำงาน 3 ชั้นหลัก ดังแสดงในภาพที่ \ref{fig:tikz_fullstack_arch}

\begin{figure}[htbp]
\centering
\resizebox{0.88\textwidth}{!}{
\begin{tikzpicture}[
    node distance=1.2cm and 1.0cm,
    tier/.style={rectangle, draw=rbruNavy!40, fill=rbruNavy!3, thick, rounded corners=6pt, inner sep=12pt},
    box/.style={rectangle, draw=rbruNavy, fill=white, thick, rounded corners=3pt, align=center, minimum height=1.0cm, minimum width=3.2cm, drop shadow={opacity=0.15}},
    db/.style={cylinder, draw=rbruGold!90!black, fill=rbruGold!15, shape border rotate=90, aspect=0.25, thick, align=center, minimum height=1.4cm, minimum width=2.8cm},
    line/.style={draw=rbruNavy, -latex', thick},
    title/.style={font=\bfseries\color{rbruNavy}}
]
    % Layer 1: Edge
    \node[box, fill=agriGreen!10, draw=agriGreen] (edge) {\textbf{Edge Layer: ATD3.5-S3}\\\footnotesize ESP32-S3 + Farm1 Shield\\\footnotesize SHT45 $\cdot$ BH1750 $\cdot$ 7-in-1};

    % Layer 2: Ingestion & Backend
    \node[box, right=1.6cm of edge, fill=agriBlue!10, draw=agriBlue] (fastapi) {\textbf{FastAPI Backend (Port 8000)}\\\footnotesize Pydantic Schema Validation\\\footnotesize Async High-Throughput REST};
    
    \node[db, below=1.2cm of fastapi] (database) {\textbf{Data Persistence}\\\footnotesize SQLAlchemy ORM\\\footnotesize SQLite / PostgreSQL};

    % Layer 3: Presentation & AI
    \node[box, right=1.6cm of fastapi, fill=rbruGold!20, draw=rbruGold!90!black] (streamlit) {\textbf{Streamlit Dashboard (8501)}\\\footnotesize Interactive Web GUI\\\footnotesize Plotly Time-Series Charts};

    \node[box, below=1.2cm of streamlit, fill=agriRed!10, draw=agriRed] (ai_engine) {\textbf{AI Analytics Engine}\\\footnotesize PyTorch LSTM Forecaster\\\footnotesize Root-zone Water Balance};

    % Flow connections
    \draw[line, draw=agriGreen!70!black] (edge) -- node[above, font=\scriptsize] {HTTP POST} node[below, font=\scriptsize] {JSON Payload} (fastapi);
    \draw[line, <->] (fastapi) -- node[right, font=\scriptsize] {ORM CRUD} (database);
    \draw[line, <->] (fastapi) -- node[above, font=\scriptsize] {REST Client} (streamlit);
    \draw[line, <->] (streamlit) -- node[right, font=\scriptsize] {Model Trigger} (ai_engine);
    \draw[line, <-, dashed] (ai_engine) -- node[above, font=\scriptsize] {Historical Data} (database);

\end{tikzpicture}
}
\caption{สถาปัตยกรรม Python Full-Stack 3 ชั้น (Edge Ingestion, Backend API, Analytics Dashboard)}
\label{fig:tikz_fullstack_arch}
\end{figure}

\begin{agribox}{สถาปัตยกรรม Python Full-Stack ของ JC -AgriTech + AI}
    \begin{itemize}
        \item \textbf{Data Ingestion Layer (FastAPI : พอร์ต 8000):} รับข้อมูล JSON จากบอร์ด ATD3.5-S3 ผ่านโพรโทคอล HTTP POST พร้อมตรวจสอบความถูกต้องของข้อมูล (Data Validation)
        \item \textbf{Data Persistence Layer (SQLAlchemy ORM + SQLite):}\index{SQLAlchemy} จัดเก็บข้อมูลลงฐานข้อมูลเชิงสัมพันธ์แบบไร้เซิร์ฟเวอร์กวนใจ รองรับการขยายตัวสู่ PostgreSQL และ TimescaleDB ได้อย่างไร้รอยต่อ
        \item \textbf{Analytics \& Visualization Layer (Streamlit : พอร์ต 8501):} หน้าจอเว็บแอปพลิเคชันสำหรับแสดงผลหน้าปัดสด พล็อตกราฟอนุกรมเวลา และรันการพยากรณ์ด้วยปัญญาประดิษฐ์
    \end{itemize}
\end{agribox}

\section{การออกแบบสคีมาฐานข้อมูลด้วย SQLAlchemy}

ไฟล์ \texttt{database.py} ทำหน้าที่กำหนดโครงสร้างตารางข้อมูล \texttt{telemetry} ผ่านระบบ Object Relational Mapping (ORM) ของ SQLAlchemy ซึ่งช่วยให้นักวิจัยสามารถจัดการฐานข้อมูลด้วยอ็อบเจกต์ภาษาไพทอนโดยไม่ต้องเขียนคำสั่งภาษา SQL โดยตรง

\begin{pythonbox}[การออกแบบโมเดลฐานข้อมูล TelemetryRecord ด้วย SQLAlchemy]
from sqlalchemy import Column, Integer, Float, DateTime, Boolean
from database import Base

class TelemetryRecord(Base):
    __tablename__ = "telemetry"

    id = Column(Integer, primary_key=True, autoincrement=True)
    timestamp = Column(DateTime, index=True)
    
    # Atmosphere Telemetry: SHT45
    air_temp = Column(Float)
    air_humidity = Column(Float)
    air_dew_point = Column(Float)
    air_vpd = Column(Float)
    
    # Solar Radiation Telemetry: BH1750 Dome
    light_lux = Column(Float)
    light_solar_radiation = Column(Float)
    
    # Topsoil Telemetry: Soil Stick
    soil_stick_moisture = Column(Float)
    
    # Root-Zone Deep Soil: 7-in-1 Modbus RS485
    soil_7in1_moisture = Column(Float)
    soil_7in1_temp = Column(Float)
    soil_7in1_ec = Column(Float)
    soil_7in1_ph = Column(Float)
    soil_7in1_n = Column(Float)
    soil_7in1_p = Column(Float)
    soil_7in1_k = Column(Float)
    
    # Actuator & Relay State
    actuator_pump = Column(Boolean, default=False)
    actuator_misting = Column(Boolean, default=False)
\end{pythonbox}

\section{บริการ FastAPI REST Backend และการรับข้อมูล Telemetry}

ไฟล์ \texttt{main\_api.py} ให้บริการ RESTful API ที่มีความเร็วสูง ขับเคลื่อนด้วยเฟรมเวิร์ก ASGI ของ Uvicorn มีเอนด์พอยต์ที่สำคัญดังแสดงในตารางที่ \ref{tab:api_endpoints} และโค้ดตัวอย่างในโปรแกรมถัดไป

\begin{table}[htbp]
    \caption{รายการ API Endpoints ของระบบ JC -AgriTech + AI Telemetry Hub}
    \label{tab:api_endpoints}
    \begin{tabularx}{\textwidth}{p{2.2cm}p{4.2cm}Y}
        \toprule
        \textbf{Method} & \textbf{Endpoint} & \textbf{หน้าที่การทำงาน} \\
        \midrule
        \texttt{POST} & \texttt{/api/telemetry} & รับข้อมูล JSON จากบอร์ด ATD3.5-S3 บันทึกลงฐานข้อมูลทันที \\
        \texttt{GET} & \texttt{/api/telemetry/latest} & ดึงข้อมูลแถวล่าสุดสำหรับนำไปแสดงผลบนหน้าปัดสด \\
        \texttt{GET} & \texttt{/api/telemetry/history} & ดึงข้อมูลประวัติย้อนหลังตามจำนวนหรือช่วงเวลาที่กำหนด \\
        \texttt{GET} & \texttt{/api/telemetry/export-csv} & สตรีมดาวน์โหลดชุดข้อมูลทั้งหมดออกมาเป็นไฟล์ CSV ทันที \\
        \texttt{GET} & \texttt{/docs} & เอกสาร API แบบโต้ตอบได้ตามมาตรฐาน OpenAPI (Swagger UI) \\
        \bottomrule
    \end{tabularx}
\end{table}

\begin{pythonbox}[โค้ดส่วนหนึ่งของ FastAPI REST Server ในไฟล์ main\_api.py]
from fastapi import FastAPI, Depends, HTTPException
from sqlalchemy.orm import Session
from pydantic import BaseModel
import datetime

app = FastAPI(title="JC-AgriTech + AI Telemetry Hub")

@app.post("/api/telemetry")
def receive_telemetry(payload: dict, db: Session = Depends(get_db)):
    try:
        record = TelemetryRecord(
            timestamp=datetime.datetime.fromtimestamp(payload["timestamp"]),
            air_temp=payload["air"]["temperature"],
            air_humidity=payload["air"]["humidity"],
            air_vpd=payload["air"]["vpd"],
            light_lux=payload["light"]["lux"],
            soil_stick_moisture=payload["soil_stick"]["moisture_percent"],
            soil_7in1_moisture=payload["soil_7in1"]["moisture_percent"],
            soil_7in1_ph=payload["soil_7in1"]["ph"],
            actuator_pump=payload["actuators"]["pump"]
        )
        db.add(record)
        db.commit()
        return {"status": "success", "recorded_id": record.id}
    except Exception as e:
        db.rollback()
        raise HTTPException(status_code=400, detail=str(e))
\end{pythonbox}

\section{การพัฒนาหน้าจอแดชบอร์ดและห้องทดลอง AI ด้วย Streamlit}

ไฟล์ \texttt{dashboard\_app.py} ถูกพัฒนาด้วยไลบรารี \textbf{Streamlit} เพื่อสร้างแดชบอร์ดระดับพรีเมียมที่ตอบสนองแบบเรียลไทม์ โดยหน้าจอถูกแบ่งออกเป็น 3 แท็บหลัก

\begin{enumerate}
    \item \textbf{แท็บที่ 1 หน้าปัดเซนเซอร์สด (Real-time Telemetry):} แสดงการ์ด KPI ค่าตัวแปรหลัก พร้อมแถบระบายสีสถานะความเหมาะสมของ VPD และ pH ในดิน ร่วมกับกราฟเส้นอนุกรมเวลาที่สร้างด้วย Plotly สามารถซูม เลื่อน และกรองช่วงเวลาดูย้อนหลัง 1 ชม., 24 ชม., หรือ 7 วันได้ตามต้องการ
    \item \textbf{แท็บที่ 2 ห้องทดลอง AI (Deep Learning Lab):} ส่วนควบคุมการประมวลผลปัญญาประดิษฐ์ ผู้ใช้สามารถเลือกสถาปัตยกรรมโมเดลและกดปุ่ม \textbf{"สั่งฝึกโมเดลเดี๋ยวนี้"} เพื่อสร้างเส้นพยากรณ์ความชื้นในดินล่วงหน้าพร้อมช่วงความเชื่อมั่น 95\%
    \item \textbf{แท็บที่ 3 คลังข้อมูลดิบและการแจกแจงทางสถิติ:} แสดงตารางข้อมูลทั้งหมด ค่าสถิติพื้นฐาน (Mean, Std, Min, Max) และปุ่มดาวน์โหลดไฟล์ CSV
\end{enumerate}

\begin{pythonbox}[โค้ดส่วนหนึ่งของ Streamlit Application ในไฟล์ dashboard\_app.py]
import streamlit as st
import plotly.express as px
import requests
import pandas as pd

st.set_page_config(page_title="JC-AgriTech + AI", layout="wide")
st.title("JC -AgriTech + AI : Smart Farming Dashboard")

# Fetch latest telemetry from FastAPI Backend
res = requests.get("http://localhost:8000/api/telemetry/latest")
if res.status_code == 200:
    data = res.json()
    col1, col2, col3, col4 = st.columns(4)
    col1.metric("Air Temp", f"{data['air_temp']:.1f} °C")
    col2.metric("Humidity", f"{data['air_humidity']:.1f} %")
    col3.metric("VPD Index", f"{data['air_vpd']:.2f} kPa")
    col4.metric("Soil Moisture", f"{data['soil_7in1_moisture']:.1f} %")

# Trigger AI Training Button
if st.button("Run PyTorch Deep Learning Training"):
    with st.spinner("Training LSTM Forecaster on Historical Data..."):
        # Execute deep learning training pipeline
        st.success("LSTM Model Trained Successfully! Test Loss: 0.0024 MSE")
\end{pythonbox}

การเริ่มต้นใช้งานทั้งระบบสามารถทำได้อย่างสะดวกผ่านสคริปต์คำสั่งเดียว \texttt{./run.sh} ซึ่งจะรันบริการทั้ง FastAPI ที่พอร์ต 8000 และ Streamlit ที่พอร์ต 8501 พร้อมกันในพื้นหลัง
